\documentclass[11pt]{article}

\usepackage[margin=1in]{geometry}
\usepackage{amsmath, amssymb, amsthm}
\usepackage{hyperref}

\newtheorem{theorem}{Theorem}
\newtheorem{proposition}[theorem]{Proposition}
\newtheorem{lemma}[theorem]{Lemma}
\theoremstyle{definition}
\newtheorem{definition}[theorem]{Definition}
\newtheorem{example}[theorem]{Example}
\theoremstyle{remark}
\newtheorem{remark}[theorem]{Remark}

\title{Derivatives (Univariate): A Rigorous Introduction}
\author{Math Foundations Map --- Node 12}
\date{}

\begin{document}
\maketitle

\section{Motivation}

Given a function $f : I \to \mathbb{R}$ on an open interval $I$ and a point
$c \in I$, we wish to assign a number $f'(c)$ that captures the
\emph{instantaneous rate of change} of $f$ at $c$. Geometrically, this number
is the slope of the line that best approximates the graph of $f$ in an
arbitrarily small neighbourhood of $(c, f(c))$. Algebraically, it is the limit
of the slopes of secant lines through $(c, f(c))$ as the second point
collapses onto $(c, f(c))$.

\section{Definitions}

\begin{definition}[Derivative at a point]
Let $f : I \to \mathbb{R}$ where $I \subseteq \mathbb{R}$ is an open interval,
and let $c \in I$. We say $f$ is \textbf{differentiable at $c$} if the limit
\[
f'(c) \;:=\; \lim_{h \to 0} \frac{f(c+h) - f(c)}{h}
\]
exists in $\mathbb{R}$. The value $f'(c)$ is called the \textbf{derivative}
of $f$ at $c$. An equivalent formulation, obtained by substituting $x = c + h$, is
\[
f'(c) \;=\; \lim_{x \to c} \frac{f(x) - f(c)}{x - c}.
\]
\end{definition}

\begin{definition}[Differentiability on an interval]
$f$ is \textbf{differentiable on $I$} if $f'(c)$ exists for every $c \in I$.
In this case the assignment $c \mapsto f'(c)$ defines a function
$f' : I \to \mathbb{R}$, called the \textbf{derivative function}.
\end{definition}

\begin{definition}[Local linear approximation]
$f$ is differentiable at $c$ with derivative $m$ if and only if there exists
a function $\varepsilon$ defined on a neighbourhood of $0$ with
$\lim_{h \to 0} \varepsilon(h) = 0$ and
\[
f(c + h) \;=\; f(c) + m \cdot h + h \cdot \varepsilon(h).
\]
\end{definition}

\section{Differentiability implies continuity}

The first nontrivial fact about differentiation is that it strengthens
continuity. We isolate this as the main theorem of the lesson.

\begin{theorem}[Differentiable implies continuous]
\label{thm:diff-implies-cts}
If $f : I \to \mathbb{R}$ is differentiable at $c \in I$, then $f$ is
continuous at $c$.
\end{theorem}

\begin{proof}
We must show $\lim_{x \to c} f(x) = f(c)$, equivalently
$\lim_{x \to c} (f(x) - f(c)) = 0$. For $x \in I$ with $x \neq c$, we use the
algebraic identity
\begin{equation}
\label{eq:factorization}
f(x) - f(c) \;=\; \frac{f(x) - f(c)}{x - c} \cdot (x - c).
\end{equation}
By the definition of $f'(c)$,
\[
\lim_{x \to c} \frac{f(x) - f(c)}{x - c} \;=\; f'(c) \;\in\; \mathbb{R}.
\]
Also, trivially,
\[
\lim_{x \to c} (x - c) \;=\; 0.
\]
The product rule for limits (which holds whenever both factors have finite
limits) applied to \eqref{eq:factorization} yields
\[
\lim_{x \to c} \bigl(f(x) - f(c)\bigr)
\;=\; \left( \lim_{x \to c} \frac{f(x) - f(c)}{x - c} \right)
      \cdot \left( \lim_{x \to c} (x - c) \right)
\;=\; f'(c) \cdot 0 \;=\; 0.
\]
Therefore $\lim_{x \to c} f(x) = f(c)$, i.e.\ $f$ is continuous at $c$.
\end{proof}

\begin{remark}
The converse is false. The function $f(x) = |x|$ is continuous at $0$ but the
two one-sided limits of the difference quotient are $+1$ and $-1$, so $f'(0)$
does not exist.
\end{remark}

\section{Differentiation rules}

\begin{proposition}[Algebraic rules]
Let $f, g$ be differentiable at $c$, and let $\alpha \in \mathbb{R}$. Then:
\begin{enumerate}
\item $(\alpha f)'(c) = \alpha f'(c)$;
\item $(f + g)'(c) = f'(c) + g'(c)$;
\item $(f g)'(c) = f'(c) g(c) + f(c) g'(c)$;
\item if $g(c) \neq 0$,
$\displaystyle (f/g)'(c) = \frac{f'(c) g(c) - f(c) g'(c)}{g(c)^2}$.
\end{enumerate}
\end{proposition}

\begin{proposition}[Chain rule]
If $g$ is differentiable at $c$ and $f$ is differentiable at $g(c)$, then
$f \circ g$ is differentiable at $c$ and
$(f \circ g)'(c) = f'(g(c)) \cdot g'(c)$.
\end{proposition}

\section{The Mean Value Theorem}

\begin{theorem}[Rolle]
If $f : [a, b] \to \mathbb{R}$ is continuous on $[a, b]$, differentiable on
$(a, b)$, and $f(a) = f(b)$, then there exists $\xi \in (a, b)$ with
$f'(\xi) = 0$.
\end{theorem}

\begin{theorem}[Mean Value Theorem]
If $f : [a, b] \to \mathbb{R}$ is continuous on $[a, b]$ and differentiable on
$(a, b)$, then there exists $\xi \in (a, b)$ with
\[
f'(\xi) \;=\; \frac{f(b) - f(a)}{b - a}.
\]
\end{theorem}

The MVT is the bridge between local information ($f'$) and global behaviour
($f(b) - f(a)$). Consequences include: monotonicity criteria, $L^\infty$
bounds on increments, Taylor's theorem with Lagrange remainder, and the
uniqueness of antiderivatives up to additive constants.

\section{Example: $\frac{d}{dx} x^2$ from first principles}

\begin{example}
For $f(x) = x^2$ and $c \in \mathbb{R}$,
\[
\frac{f(c+h) - f(c)}{h}
= \frac{(c+h)^2 - c^2}{h}
= \frac{2ch + h^2}{h}
= 2c + h
\;\xrightarrow{h \to 0}\; 2c.
\]
Hence $f'(c) = 2c$, i.e.\ $\tfrac{d}{dx} x^2 = 2x$.
\end{example}

\end{document}
